\documentclass[11pt]{standalone}

\usepackage{ifthen}
\usepackage{amsmath}
\usepackage{tikz} 
\usepackage[d]{esvect}
\usetikzlibrary{shapes.misc}
\usetikzlibrary{arrows,arrows.meta}
\usetikzlibrary{calc,intersections, patterns, math}

\definecolor{pfeil}{RGB}{168,167,167}
\definecolor{salmon}{RGB}{250,128,114}
\definecolor{petrol}{RGB}{0, 118, 136}
\definecolor{darkgoldenrod}{RGB}{184, 134, 11}
\colorlet{petrol-lighter}{petrol!40}
\colorlet{darkgoldenrod-lighter}{darkgoldenrod!40}

\begin{document}

\begin{tikzpicture}[pfeil, >=stealth, very thick]
			
	\draw[magenta,->] (1,1) -- node[below] {$\vec b$} ++ (3,1);

	\draw[magenta,dashed,ultra thick,->] (1,1) ++(3,1) -- node[below] {$\vec b$} ++ (3,1); 
	\draw[magenta,dashed,ultra thick,->] (1,1) ++(3,1)++(3,1) -- node[below] {$\vec b$} ++ (3,1);
	% \draw[magenta,dashed,ultra thick,->] (1,1) ++(3,1) -- node[below] {$\vec b$} ++ (3,1); 
	% \draw[magenta,dashed,ultra thick,->] (1,1) ++(3,1) ++(3,1) -- node[below] {$\vec ab$} ++ (3,1); 
			
	\draw[green,thick,->] (1,1) -- node[above, sloped] {$\vec c = {\color{orange}a}\cdot {\color{magenta}\vec b}$} (10,4);
\end{tikzpicture}			

\end{document}
